% --- Archon physics formalization source begin ---
% archon:physics
% archon:covers PhyXMiniProblems/problem_phyx_mini_0892.lean
% archon:source-report reports/phyx_mini/problem_phyx_mini_0892.source.json

\chapter{Physics problem phyx\_mini\_0892}
\label{ch:PhyXMiniProblems_problem_phyx_mini_0892}

\paragraph{Problem source.}
\#\# Question

What is the nit electric force on charge A?

\#\# Answer choices

- A: A: 1.25N

- B: B: 1.35N

- C: C: 0N

- D: D: 5.33N

\#\# Figure caption (auxiliary; use the image as primary evidence)

The image shows three charged objects labeled A, B, and C arranged in a straight line. 

- Object A is red with a plus sign inside, indicating a positive charge. It is labeled as having a charge of \textbackslash{}(1.0 \textbackslash{}, \textbackslash{}text\{nC\}\textbackslash{}) (nanocoulombs).
- Object B is teal with a minus sign inside, indicating a negative charge. It is labeled as having a charge of \textbackslash{}(-1.0 \textbackslash{}, \textbackslash{}text\{nC\}\textbackslash{}).
- Object C is also red with a plus sign inside, indicating a positive charge. It is labeled as having a charge of \textbackslash{}(4.0 \textbackslash{}, \textbackslash{}text\{nC\}\textbackslash{}).

The distance between A and B is shown as \textbackslash{}(1.0 \textbackslash{}, \textbackslash{}text\{cm\}\textbackslash{}), and the distance between B and C is also \textbackslash{}(1.0 \textbackslash{}, \textbackslash{}text\{cm\}\textbackslash{}). Arrows are used to indicate these distances.

\#\# Dataset metadata

Electrostatics; Spatial Relation Reasoning, Multi-Formula Reasoning

\paragraph{Recorded answer/context.}
C: C: 0N

\paragraph{Figure/image path.}
/root/proposal\_for\_physic/science-mango/phyx\_mini\_run/phyx\_data/test\_image/892.png

\paragraph{Formalization target.}
create a compiling Lean file with sorry bodies at `PhyXMiniProblems/problem\_phyx\_mini\_0892.lean`. The Lean declarations must preserve the physical quantities, dimensions or dimensional roles, figure labels, governing-law hypotheses, and final relation expressed by this problem.
Use Mathlib/Physlib names found through LeanExplore where available. If a physics API is missing, introduce faithful local abstractions rather than scalar placeholder aliases.

\begin{theorem}[Physics formalization target]
\label{thm:physics:phyx_mini_0892:target}
\lean{PhyXMiniProblems.ProblemPhyXMini0892.problem_phyx_mini_0892}
\uses{def:physics:phyx-mini-0892:phyxminiproblems-problemphyxmini0892-forcecomponentinnewtons, def:physics:phyx-mini-0892:phyxminiproblems-problemphyxmini0892-threepointchargelinesetup, def:physics:phyx-mini-0892:phyxminiproblems-problemphyxmini0892-matchesthreepointchargescenario, def:physics:phyx-mini-0892:phyxminiproblems-problemphyxmini0892-matchessuppliedthreechargefigure, def:physics:phyx-mini-0892:phyxminiproblems-problemphyxmini0892-hascollinearadjacentgapgeometry, def:physics:phyx-mini-0892:phyxminiproblems-problemphyxmini0892-hasphysicalelectrostaticparameters, def:physics:phyx-mini-0892:phyxminiproblems-problemphyxmini0892-satisfiesaxialpointchargeforcelaw, def:physics:phyx-mini-0892:phyxminiproblems-problemphyxmini0892-satisfiesaxialforcesuperposition}
The assigned autoformalize agent should translate this physics problem into Lean declarations in the covered file, with theorem and lemma proof bodies written as `by sorry`.
\end{theorem}
\begin{proof}
This is an autoformalization task, not a proof task. Produce statements that can later be proved by the physics prover without weakening the physical model.
\end{proof}

% --- Archon declaration topology begin ---
\section{Declaration topology}
\begin{definition}[force Dimension]
  \label{def:physics:phyx-mini-0892:phyxminiproblems-problemphyxmini0892-forcedimension}
  \lean{PhyXMiniProblems.ProblemPhyXMini0892.forceDimension}
  The physical dimension M L T⁻² of force
\end{definition}
\begin{proof}
  This is the declared model definition of force Dimension.
\end{proof}
\begin{definition}[Signed Charge Quantity]
  \label{def:physics:phyx-mini-0892:phyxminiproblems-problemphyxmini0892-signedchargequantity}
  \lean{PhyXMiniProblems.ProblemPhyXMini0892.SignedChargeQuantity}
  A signed, unit-independent electric charge
\end{definition}
\begin{proof}
  This is the declared model definition of Signed Charge Quantity.
\end{proof}
\begin{definition}[Axial Position Quantity]
  \label{def:physics:phyx-mini-0892:phyxminiproblems-problemphyxmini0892-axialpositionquantity}
  \lean{PhyXMiniProblems.ProblemPhyXMini0892.AxialPositionQuantity}
  A signed position coordinate on the line containing the charges
\end{definition}
\begin{proof}
  This is the declared model definition of Axial Position Quantity.
\end{proof}
\begin{definition}[Length Quantity]
  \label{def:physics:phyx-mini-0892:phyxminiproblems-problemphyxmini0892-lengthquantity}
  \lean{PhyXMiniProblems.ProblemPhyXMini0892.LengthQuantity}
  A nonnegative, unit-independent physical separation
\end{definition}
\begin{proof}
  This is the declared model definition of Length Quantity.
\end{proof}
\begin{definition}[Axial Force Component Quantity]
  \label{def:physics:phyx-mini-0892:phyxminiproblems-problemphyxmini0892-axialforcecomponentquantity}
  \lean{PhyXMiniProblems.ProblemPhyXMini0892.AxialForceComponentQuantity}
  \uses{def:physics:phyx-mini-0892:phyxminiproblems-problemphyxmini0892-forcedimension}
  A signed axial component of a physical force
\end{definition}
\begin{proof}
  This is the declared model definition of Axial Force Component Quantity.
\end{proof}
\begin{definition}[charge In Coulombs]
  \label{def:physics:phyx-mini-0892:phyxminiproblems-problemphyxmini0892-chargeincoulombs}
  \lean{PhyXMiniProblems.ProblemPhyXMini0892.chargeInCoulombs}
  \uses{def:physics:phyx-mini-0892:phyxminiproblems-problemphyxmini0892-signedchargequantity}
  Read a signed physical charge in coherent-SI coulombs
\end{definition}
\begin{proof}
  This is the declared model definition of charge In Coulombs.
\end{proof}
\begin{definition}[charge In Nanocoulombs]
  \label{def:physics:phyx-mini-0892:phyxminiproblems-problemphyxmini0892-chargeinnanocoulombs}
  \lean{PhyXMiniProblems.ProblemPhyXMini0892.chargeInNanocoulombs}
  \uses{def:physics:phyx-mini-0892:phyxminiproblems-problemphyxmini0892-signedchargequantity, def:physics:phyx-mini-0892:phyxminiproblems-problemphyxmini0892-chargeincoulombs}
  Read a charge in the nanocoulombs used by the figure labels
\end{definition}
\begin{proof}
  This is the declared model definition of charge In Nanocoulombs.
\end{proof}
\begin{definition}[position In Meters]
  \label{def:physics:phyx-mini-0892:phyxminiproblems-problemphyxmini0892-positioninmeters}
  \lean{PhyXMiniProblems.ProblemPhyXMini0892.positionInMeters}
  \uses{def:physics:phyx-mini-0892:phyxminiproblems-problemphyxmini0892-axialpositionquantity}
  Read an axial position in coherent-SI metres
\end{definition}
\begin{proof}
  This is the declared model definition of position In Meters.
\end{proof}
\begin{definition}[length In Meters]
  \label{def:physics:phyx-mini-0892:phyxminiproblems-problemphyxmini0892-lengthinmeters}
  \lean{PhyXMiniProblems.ProblemPhyXMini0892.lengthInMeters}
  \uses{def:physics:phyx-mini-0892:phyxminiproblems-problemphyxmini0892-lengthquantity}
  Read a nonnegative physical length in coherent-SI metres
\end{definition}
\begin{proof}
  This is the declared model definition of length In Meters.
\end{proof}
\begin{definition}[length In Centimeters]
  \label{def:physics:phyx-mini-0892:phyxminiproblems-problemphyxmini0892-lengthincentimeters}
  \lean{PhyXMiniProblems.ProblemPhyXMini0892.lengthInCentimeters}
  \uses{def:physics:phyx-mini-0892:phyxminiproblems-problemphyxmini0892-lengthquantity, def:physics:phyx-mini-0892:phyxminiproblems-problemphyxmini0892-lengthinmeters}
  Read a length in the centimetres used by the distance arrows
\end{definition}
\begin{proof}
  This is the declared model definition of length In Centimeters.
\end{proof}
\begin{definition}[force Component In Newtons]
  \label{def:physics:phyx-mini-0892:phyxminiproblems-problemphyxmini0892-forcecomponentinnewtons}
  \lean{PhyXMiniProblems.ProblemPhyXMini0892.forceComponentInNewtons}
  \uses{def:physics:phyx-mini-0892:phyxminiproblems-problemphyxmini0892-axialforcecomponentquantity}
  Read a signed axial force component in coherent-SI newtons
\end{definition}
\begin{proof}
  This is the declared model definition of force Component In Newtons.
\end{proof}
\begin{definition}[Charge Site]
  \label{def:physics:phyx-mini-0892:phyxminiproblems-problemphyxmini0892-chargesite}
  \lean{PhyXMiniProblems.ProblemPhyXMini0892.ChargeSite}
  The three charged objects, named exactly as in image 892.png
\end{definition}
\begin{proof}
  This is the declared model definition of Charge Site.
\end{proof}
\begin{definition}[Source For A]
  \label{def:physics:phyx-mini-0892:phyxminiproblems-problemphyxmini0892-sourcefora}
  \lean{PhyXMiniProblems.ProblemPhyXMini0892.SourceForA}
  The two charges other than A which exert force on A
\end{definition}
\begin{proof}
  This is the declared model definition of Source For A.
\end{proof}
\begin{definition}[to Charge Site]
  \label{def:physics:phyx-mini-0892:phyxminiproblems-problemphyxmini0892-sourcefora-tochargesite}
  \lean{PhyXMiniProblems.ProblemPhyXMini0892.SourceForA.toChargeSite}
  \uses{def:physics:phyx-mini-0892:phyxminiproblems-problemphyxmini0892-chargesite, def:physics:phyx-mini-0892:phyxminiproblems-problemphyxmini0892-sourcefora}
  Regard a source of force on A as one of the three labelled sites
\end{definition}
\begin{proof}
  This is the declared model definition of to Charge Site.
\end{proof}
\begin{definition}[Adjacent Gap]
  \label{def:physics:phyx-mini-0892:phyxminiproblems-problemphyxmini0892-adjacentgap}
  \lean{PhyXMiniProblems.ProblemPhyXMini0892.AdjacentGap}
  The two adjacent gaps whose arrows are printed in the image
\end{definition}
\begin{proof}
  This is the declared model definition of Adjacent Gap.
\end{proof}
\begin{definition}[left Endpoint]
  \label{def:physics:phyx-mini-0892:phyxminiproblems-problemphyxmini0892-adjacentgap-leftendpoint}
  \lean{PhyXMiniProblems.ProblemPhyXMini0892.AdjacentGap.leftEndpoint}
  \uses{def:physics:phyx-mini-0892:phyxminiproblems-problemphyxmini0892-chargesite, def:physics:phyx-mini-0892:phyxminiproblems-problemphyxmini0892-adjacentgap}
  Left endpoint of a displayed adjacent gap
\end{definition}
\begin{proof}
  This is the declared model definition of left Endpoint.
\end{proof}
\begin{definition}[right Endpoint]
  \label{def:physics:phyx-mini-0892:phyxminiproblems-problemphyxmini0892-adjacentgap-rightendpoint}
  \lean{PhyXMiniProblems.ProblemPhyXMini0892.AdjacentGap.rightEndpoint}
  \uses{def:physics:phyx-mini-0892:phyxminiproblems-problemphyxmini0892-chargesite, def:physics:phyx-mini-0892:phyxminiproblems-problemphyxmini0892-adjacentgap}
  Right endpoint of a displayed adjacent gap
\end{definition}
\begin{proof}
  This is the declared model definition of right Endpoint.
\end{proof}
\begin{definition}[Figure Charge Sign]
  \label{def:physics:phyx-mini-0892:phyxminiproblems-problemphyxmini0892-figurechargesign}
  \lean{PhyXMiniProblems.ProblemPhyXMini0892.FigureChargeSign}
  The sign glyph drawn inside a charge circle
\end{definition}
\begin{proof}
  This is the declared model definition of Figure Charge Sign.
\end{proof}
\begin{definition}[Figure Charge Color]
  \label{def:physics:phyx-mini-0892:phyxminiproblems-problemphyxmini0892-figurechargecolor}
  \lean{PhyXMiniProblems.ProblemPhyXMini0892.FigureChargeColor}
  The two circle colors visibly used to distinguish the charges
\end{definition}
\begin{proof}
  This is the declared model definition of Figure Charge Color.
\end{proof}
\begin{definition}[expected Site Label]
  \label{def:physics:phyx-mini-0892:phyxminiproblems-problemphyxmini0892-expectedsitelabel}
  \lean{PhyXMiniProblems.ProblemPhyXMini0892.expectedSiteLabel}
  \uses{def:physics:phyx-mini-0892:phyxminiproblems-problemphyxmini0892-chargesite}
  The literal label expected beside each charged object
\end{definition}
\begin{proof}
  This is the declared model definition of expected Site Label.
\end{proof}
\begin{definition}[expected Charge Sign]
  \label{def:physics:phyx-mini-0892:phyxminiproblems-problemphyxmini0892-expectedchargesign}
  \lean{PhyXMiniProblems.ProblemPhyXMini0892.expectedChargeSign}
  \uses{def:physics:phyx-mini-0892:phyxminiproblems-problemphyxmini0892-chargesite, def:physics:phyx-mini-0892:phyxminiproblems-problemphyxmini0892-figurechargesign}
  Charge signs transcribed from the primary image
\end{definition}
\begin{proof}
  This is the declared model definition of expected Charge Sign.
\end{proof}
\begin{definition}[expected Charge Color]
  \label{def:physics:phyx-mini-0892:phyxminiproblems-problemphyxmini0892-expectedchargecolor}
  \lean{PhyXMiniProblems.ProblemPhyXMini0892.expectedChargeColor}
  \uses{def:physics:phyx-mini-0892:phyxminiproblems-problemphyxmini0892-chargesite, def:physics:phyx-mini-0892:phyxminiproblems-problemphyxmini0892-figurechargecolor}
  Circle colors transcribed from the primary image
\end{definition}
\begin{proof}
  This is the declared model definition of expected Charge Color.
\end{proof}
\begin{definition}[expected Charge Label Nanocoulombs]
  \label{def:physics:phyx-mini-0892:phyxminiproblems-problemphyxmini0892-expectedchargelabelnanocoulombs}
  \lean{PhyXMiniProblems.ProblemPhyXMini0892.expectedChargeLabelNanocoulombs}
  \uses{def:physics:phyx-mini-0892:phyxminiproblems-problemphyxmini0892-chargesite}
  Nanocoulomb values printed above the three charge circles
\end{definition}
\begin{proof}
  This is the declared model definition of expected Charge Label Nanocoulombs.
\end{proof}
\begin{definition}[Three Charge Line Figure]
  \label{def:physics:phyx-mini-0892:phyxminiproblems-problemphyxmini0892-threechargelinefigure}
  \lean{PhyXMiniProblems.ProblemPhyXMini0892.ThreeChargeLineFigure}
  \uses{def:physics:phyx-mini-0892:phyxminiproblems-problemphyxmini0892-chargesite, def:physics:phyx-mini-0892:phyxminiproblems-problemphyxmini0892-adjacentgap, def:physics:phyx-mini-0892:phyxminiproblems-problemphyxmini0892-figurechargesign, def:physics:phyx-mini-0892:phyxminiproblems-problemphyxmini0892-figurechargecolor}
  Literal presentation data transcribed from image 892.png
\end{definition}
\begin{proof}
  This is the declared model definition of Three Charge Line Figure.
\end{proof}
\begin{definition}[Three Point Charge Line Setup]
  \label{def:physics:phyx-mini-0892:phyxminiproblems-problemphyxmini0892-threepointchargelinesetup}
  \lean{PhyXMiniProblems.ProblemPhyXMini0892.ThreePointChargeLineSetup}
  \uses{def:physics:phyx-mini-0892:phyxminiproblems-problemphyxmini0892-signedchargequantity, def:physics:phyx-mini-0892:phyxminiproblems-problemphyxmini0892-axialpositionquantity, def:physics:phyx-mini-0892:phyxminiproblems-problemphyxmini0892-lengthquantity, def:physics:phyx-mini-0892:phyxminiproblems-problemphyxmini0892-axialforcecomponentquantity, def:physics:phyx-mini-0892:phyxminiproblems-problemphyxmini0892-chargesite, def:physics:phyx-mini-0892:phyxminiproblems-problemphyxmini0892-sourcefora, def:physics:phyx-mini-0892:phyxminiproblems-problemphyxmini0892-adjacentgap, def:physics:phyx-mini-0892:phyxminiproblems-problemphyxmini0892-threechargelinefigure}
  The physical charges, positions, separations, and force observables associated with the diagram. The force from each other charge and the net force on A are primitive observables here; the governing laws below relate their SI readouts without defining them from the requested zero value
\end{definition}
\begin{proof}
  This is the declared model definition of Three Point Charge Line Setup.
\end{proof}
\begin{definition}[displacement From Source To AIn Meters]
  \label{def:physics:phyx-mini-0892:phyxminiproblems-problemphyxmini0892-displacementfromsourcetoainmeters}
  \lean{PhyXMiniProblems.ProblemPhyXMini0892.displacementFromSourceToAInMeters}
  \uses{def:physics:phyx-mini-0892:phyxminiproblems-problemphyxmini0892-positioninmeters, def:physics:phyx-mini-0892:phyxminiproblems-problemphyxmini0892-sourcefora, def:physics:phyx-mini-0892:phyxminiproblems-problemphyxmini0892-threepointchargelinesetup}
  Axial displacement from a source charge to charge A, in metres
\end{definition}
\begin{proof}
  This is the declared model definition of displacement From Source To AIn Meters.
\end{proof}
\begin{definition}[Matches Three Point Charge Scenario]
  \label{def:physics:phyx-mini-0892:phyxminiproblems-problemphyxmini0892-matchesthreepointchargescenario}
  \lean{PhyXMiniProblems.ProblemPhyXMini0892.MatchesThreePointChargeScenario}
  \uses{def:physics:phyx-mini-0892:phyxminiproblems-problemphyxmini0892-threepointchargelinesetup}
  All three pictured objects are modelled as point charges
\end{definition}
\begin{proof}
  This is the declared model definition of Matches Three Point Charge Scenario.
\end{proof}
\begin{definition}[Matches Supplied Three Charge Figure]
  \label{def:physics:phyx-mini-0892:phyxminiproblems-problemphyxmini0892-matchessuppliedthreechargefigure}
  \lean{PhyXMiniProblems.ProblemPhyXMini0892.MatchesSuppliedThreeChargeFigure}
  \uses{def:physics:phyx-mini-0892:phyxminiproblems-problemphyxmini0892-chargeinnanocoulombs, def:physics:phyx-mini-0892:phyxminiproblems-problemphyxmini0892-lengthincentimeters, def:physics:phyx-mini-0892:phyxminiproblems-problemphyxmini0892-expectedsitelabel, def:physics:phyx-mini-0892:phyxminiproblems-problemphyxmini0892-expectedchargesign, def:physics:phyx-mini-0892:phyxminiproblems-problemphyxmini0892-expectedchargecolor, def:physics:phyx-mini-0892:phyxminiproblems-problemphyxmini0892-expectedchargelabelnanocoulombs, def:physics:phyx-mini-0892:phyxminiproblems-problemphyxmini0892-threepointchargelinesetup}
  The literal image features and their calibration to the independent physical charges and separations. These premises contain no force value
\end{definition}
\begin{proof}
  This is the declared model definition of Matches Supplied Three Charge Figure.
\end{proof}
\begin{definition}[Has Collinear Adjacent Gap Geometry]
  \label{def:physics:phyx-mini-0892:phyxminiproblems-problemphyxmini0892-hascollinearadjacentgapgeometry}
  \lean{PhyXMiniProblems.ProblemPhyXMini0892.HasCollinearAdjacentGapGeometry}
  \uses{def:physics:phyx-mini-0892:phyxminiproblems-problemphyxmini0892-positioninmeters, def:physics:phyx-mini-0892:phyxminiproblems-problemphyxmini0892-lengthinmeters, def:physics:phyx-mini-0892:phyxminiproblems-problemphyxmini0892-threepointchargelinesetup}
  The signed position coordinates realize the printed A--B--C order and the two positive physical separations
\end{definition}
\begin{proof}
  This is the declared model definition of Has Collinear Adjacent Gap Geometry.
\end{proof}
\begin{definition}[Has Physical Electrostatic Parameters]
  \label{def:physics:phyx-mini-0892:phyxminiproblems-problemphyxmini0892-hasphysicalelectrostaticparameters}
  \lean{PhyXMiniProblems.ProblemPhyXMini0892.HasPhysicalElectrostaticParameters}
  \uses{def:physics:phyx-mini-0892:phyxminiproblems-problemphyxmini0892-chargeincoulombs, def:physics:phyx-mini-0892:phyxminiproblems-problemphyxmini0892-threepointchargelinesetup}
  Basic nondegeneracy and positivity conditions for the electrostatic model
\end{definition}
\begin{proof}
  This is the declared model definition of Has Physical Electrostatic Parameters.
\end{proof}
\begin{definition}[Satisfies Axial Point Charge Force Law]
  \label{def:physics:phyx-mini-0892:phyxminiproblems-problemphyxmini0892-satisfiesaxialpointchargeforcelaw}
  \lean{PhyXMiniProblems.ProblemPhyXMini0892.SatisfiesAxialPointChargeForceLaw}
  \uses{def:physics:phyx-mini-0892:phyxminiproblems-problemphyxmini0892-chargeincoulombs, def:physics:phyx-mini-0892:phyxminiproblems-problemphyxmini0892-positioninmeters, def:physics:phyx-mini-0892:phyxminiproblems-problemphyxmini0892-forcecomponentinnewtons, def:physics:phyx-mini-0892:phyxminiproblems-problemphyxmini0892-threepointchargelinesetup, def:physics:phyx-mini-0892:phyxminiproblems-problemphyxmini0892-displacementfromsourcetoainmeters}
  For each source S ∈ \{B,C\}, the signed one-dimensional force on A is k q\_A q\_S (x\_A - x\_S) / |x\_A - x\_S|\textasciicircum{}3. This is Coulomb's inverse-square law with the displacement factor retaining the force direction. It is a general governing law, not the target value
\end{definition}
\begin{proof}
  This is the declared model definition of Satisfies Axial Point Charge Force Law.
\end{proof}
\begin{definition}[Satisfies Axial Force Superposition]
  \label{def:physics:phyx-mini-0892:phyxminiproblems-problemphyxmini0892-satisfiesaxialforcesuperposition}
  \lean{PhyXMiniProblems.ProblemPhyXMini0892.SatisfiesAxialForceSuperposition}
  \uses{def:physics:phyx-mini-0892:phyxminiproblems-problemphyxmini0892-forcecomponentinnewtons, def:physics:phyx-mini-0892:phyxminiproblems-problemphyxmini0892-sourcefora, def:physics:phyx-mini-0892:phyxminiproblems-problemphyxmini0892-threepointchargelinesetup}
  The signed net force on A is the sum of the forces due to B and C
\end{definition}
\begin{proof}
  This is the declared model definition of Satisfies Axial Force Superposition.
\end{proof}
\begin{lemma}[source force components cancel]
  \label{lem:physics:phyx-mini-0892:phyxminiproblems-problemphyxmini0892-source-force-components-cancel}
  \lean{PhyXMiniProblems.ProblemPhyXMini0892.source_force_components_cancel}
  \uses{def:physics:phyx-mini-0892:phyxminiproblems-problemphyxmini0892-forcecomponentinnewtons, def:physics:phyx-mini-0892:phyxminiproblems-problemphyxmini0892-threepointchargelinesetup, def:physics:phyx-mini-0892:phyxminiproblems-problemphyxmini0892-matchessuppliedthreechargefigure, def:physics:phyx-mini-0892:phyxminiproblems-problemphyxmini0892-hascollinearadjacentgapgeometry, def:physics:phyx-mini-0892:phyxminiproblems-problemphyxmini0892-hasphysicalelectrostaticparameters, def:physics:phyx-mini-0892:phyxminiproblems-problemphyxmini0892-satisfiesaxialpointchargeforcelaw}
  The -1 nC charge at 1 cm attracts A to the right, while the +4 nC charge at 2 cm repels A to the left. Coulomb's inverse-square scaling makes the two signed source-force components cancel
\end{lemma}
\begin{proof}
  The conclusion follows from the governing relations and figure readouts named in the statement of source force components cancel.
\end{proof}
\begin{definition}[Answer Choice]
  \label{def:physics:phyx-mini-0892:phyxminiproblems-problemphyxmini0892-answerchoice}
  \lean{PhyXMiniProblems.ProblemPhyXMini0892.AnswerChoice}
  Labels attached to the four force choices in the problem source
\end{definition}
\begin{proof}
  This is the declared model definition of Answer Choice.
\end{proof}
\begin{definition}[force In Newtons]
  \label{def:physics:phyx-mini-0892:phyxminiproblems-problemphyxmini0892-answerchoice-forceinnewtons}
  \lean{PhyXMiniProblems.ProblemPhyXMini0892.AnswerChoice.forceInNewtons}
  \uses{def:physics:phyx-mini-0892:phyxminiproblems-problemphyxmini0892-answerchoice}
  Force magnitude in newtons printed beside each displayed answer label
\end{definition}
\begin{proof}
  This is the declared model definition of force In Newtons.
\end{proof}
\begin{definition}[recorded Dataset Answer]
  \label{def:physics:phyx-mini-0892:phyxminiproblems-problemphyxmini0892-recordeddatasetanswer}
  \lean{PhyXMiniProblems.ProblemPhyXMini0892.recordedDatasetAnswer}
  \uses{def:physics:phyx-mini-0892:phyxminiproblems-problemphyxmini0892-answerchoice}
  The answer label recorded by the supplied dataset
\end{definition}
\begin{proof}
  This is the declared model definition of recorded Dataset Answer.
\end{proof}
% --- Archon declaration topology end ---
% --- Archon physics formalization source end ---
